\section{Robust control flow, gaps, and charts}
\label{sec:control}

\subsection{Why exact trichotomy is the wrong interface}

A finite algorithm can compare rationals exactly.  It cannot in general decide
whether two arbitrary regular names are equal, nor choose the sign of a name
that may be zero.  Many classical proofs use such a choice only to select a
formula, pivot, or coordinate chart.  The constructive repair is to make the
charts overlap, or to require a positive gap when a genuinely exact decision
is part of the conclusion.

\subsection{The overlapping branch lemma}

\begin{theorem}[Overlapping rational branch]
\label{thm:overlap-branch}
Let \(x\) be represented by a rational approximation \(a\) with
\(\abs{x-a}\le e\).  Let \(q\in\Q\), \(d>0\), and suppose \(4e\le d\).
A decidable comparison of \(a\) with \(q+d/2\) produces one of the valid
conclusions
\[
 x<q+d\qquad\text{or}\qquad x>q.
\]
The conclusions overlap on \((q,q+d)\).
\end{theorem}

\begin{proof}
If \(a\le q+d/2\), then
\[
 x\le a+e\le q+d/2+d/4<q+d.
\]
If \(a>q+d/2\), then
\[
 x\ge a-e>q+d/2-d/4>q.
\]
Only a rational comparison is performed.  In a pure name formulation, replace
strict inequalities by their rational-margin certificates.  All displayed
inequalities are checked after clearing denominators.
\end{proof}

This lemma is enough for overlapping analytic charts, small/large parameter
cases, approximate sign tests, and robust pivot selection.  It avoids the
false demand that the boundary case be assigned uniquely.

\subsection{Exact decisions from apartness}

\begin{theorem}[Exact branch under a supplied gap]
\label{thm:exact-gap}
Suppose a contract supplies \(d>0\) with
\(\abs{x-q}\ge d\).  If \(a\) approximates \(x\) within \(e<d/2\), then the
rational sign of \(a-q\) is nonzero and equals the sign of \(x-q\).
\end{theorem}

\begin{proof}
If \(x\ge q+d\), then \(a>q+d-e>q\).  If
\(x\le q-d\), then \(a<q-d+e<q\).  Conversely, an approximation equal to
\(q\) would be within \(e<d\) of \(x\), contradicting the gap.  The two cases
in the gap certificate may be encoded as a signed lower-bound witness or as a
Boolean comparison already supplied by the certificate.
\end{proof}

Apartness is therefore the correct input to inversion, normalization, root
ordering, exact sign, and uniqueness theorems.  Negative information
\(x\ne q\) is not interchangeable with a positive rational gap.

\subsection{Stable finite pivoting}

Let \(v_0,\ldots,v_{m-1}\) be rationally approximated candidate vectors, and
let \(r_i\) be a nonnegative residual or squared norm.  A finite maximum scan
returns the least index \(j\) attaining the largest sampled value.

\begin{lemma}[Robust pivot]
\label{lem:robust-pivot}
Suppose some exact residual satisfies \(r_i\ge d\), every sampled residual
\(\tilde r_i\) has error at most \(e\), and \(4e\le d\).  Then the selected
maximum satisfies \(r_j\ge d/2\).  Consequently any operation requiring a
positive lower bound may use the certified bound \(d/2\).
\end{lemma}

\begin{proof}
For the good index \(i\), \(\tilde r_i\ge d-e\).  Maximality gives
\(\tilde r_j\ge d-e\).  Hence
\(r_j\ge\tilde r_j-e\ge d-2e\ge d/2\).  The maximum scan invariant says that
the stored index dominates all processed samples and is a quantifier-free
Boolean fold.  No exact comparison of represented residuals is used.
\end{proof}

This is the generic core of stable column selection in matrix arguments,
choice of a nonzero polynomial coefficient, denominator selection in projective
coordinates, and finite chart selection on compact domains.

\subsection{Overlapping atlases}

A constructive atlas should be designed so that every point has at least one
chart with a quantitative denominator bound, and the chart-selection tests
have overlap.  For example, two formulas valid respectively when
\(c>-3/4\) and \(c<3/4\) cover all \(c\in[-1,1]\).  Approximate \(c\) to
error at most \(1/16\) and compare it to zero.  Either selected chart has a
fixed denominator margin.  On the overlap both are valid, so no proof of
unique chart selection is needed.

The same design applies to:
\begin{itemize}
\item quaternion charts for rotations;
\item affine charts of projective space;
\item inverse trigonometric formulas;
\item root formulas with competing dominant coefficients;
\item contour decompositions and local power-series disks.
\end{itemize}

Transition identities on overlaps should be finite polynomial/rational
certificates.  This replaces topological choice by a bounded rational decision.

\subsection{Projective and homogeneous output}

Normalization often creates an artificial discontinuity.  A nonzero kernel
vector has a continuous projective class even when no canonical unit vector is
available.  A polynomial root can be returned as an isolating disk or cluster
before an ordering is chosen.  A rotation can be returned as a unit quaternion
pair with vector part allowed to vanish at the identity.  A line or conic can
be represented by homogeneous coefficients with a positive maximum-norm
certificate, postponing division until a chart has been selected.

\begin{principle}[Delay normalization]
Whenever a theorem is invariant under scaling, sign, permutation, or another
finite symmetry, state and prove the invariant conclusion first.  Add a
canonical representative only under a separation or orientation certificate.
\end{principle}

This principle converts many apparent logical obstacles into ordinary finite
algebra.

\subsection{Partial operations as certified transducers}

A guarded inverse transducer takes a name \(x\), a positive rational \(d\),
and a proof that \(\abs{x}\ge d\).  Its precision transformer requests enough
input accuracy that the rational approximation is at least \(d/2\) away from
zero; it then uses rational division.  The trace records the comparison and
cross-multiplied error inequality.  Square root, logarithm, matrix inversion,
and normalization have the same shape:

\[
\text{domain/apartness bit}
\;\longrightarrow\;
\text{precision request}
\;\longrightarrow\;
\text{finite rational operation}
\;\longrightarrow\;
\text{error certificate}.
\]

The operation itself remains total on all codes, but only the guarded theorem
assigns it the intended mathematical meaning.
